% SPDX-License-Identifier: Apache-2.0
% covers: LineStore
\datasheet{LineStore}{A host on the bus that writes a run of memory
from a channel, in bursts, to a count of bytes.}

\dsfacts{\code{lib/parts/src/dma.rs}}{\code{txhdl\_parts::dma::LineStore}}%
{\code{//docs:examples}}

\subsection*{Function}

The mirror of \code{LineFetch}, and the half an Ethernet receiver
wants. Words arrive on a channel and are written to memory as bursts,
so a peripheral receiving bulk data puts it away itself rather than
interrupting the core per word.

\subsection*{Parameters}

\begin{center}\footnotesize
\begin{tabular}{@{}lp{0.68\textwidth}@{}}
\toprule
Parameter & Meaning \\
\midrule
\code{A} & The width of an address, in bits. \\
\code{I} & The width of a burst identifier, in bits. \\
\code{BEATS} & How many beats a burst asks for. \\
\code{WC} & The width of the byte count. \\
\bottomrule
\end{tabular}
\end{center}

The tables below use \code{LineStore<32, 2, 16, 16>}, the engine of
\code{ex\_dmaw}.

\dstables{LineStore}

\subsection*{Behaviour}

\begin{itemize}
\item A run starts when \code{go} is high and nothing is in flight.
\code{base} and \code{bytes} are read once, there.
\item The beats wanted are \code{(bytes + 3) / 4}: a partial last
word is still a beat, it is just not a whole one.
\item A beat goes out when a word is offered and the beat channel has
room. A producer that stops offering stops the burst, which is the
backpressure in this direction.
\item The last beat of the run is strobed down to the bytes that are
real: \code{0x1}, \code{0x3} or \code{0x7}, and \code{0xf} when the
count is a whole number of words.
\end{itemize}

\subsection*{Two differences from \code{LineFetch}, both the bus}

\emph{A write burst is finished by its response, not by its last
beat.} The host answers a read on the beat channel and a write on
\code{done}, so the identifier goes back when the response arrives,
which may be well after the last beat went out.

\emph{AXI4 puts no identifier on the beats}, so a burst's beats
belong to the oldest address phase outstanding. A second burst issued
before the first one's beats were out would take them. \code{LineFetch}
keeps one burst in flight because a second buys nothing; here it is a
rule rather than a choice.

\subsection*{Verification}

\begin{itemize}
\item \code{ex\_dmaw} writes sixty-four words across a real link with
\code{AxiHost}, \code{AxiPer} and a memory on the far side, then
reads the memory back. The netlist is checked against that run under
nvc and Verilator: \code{//docs:dmaw\_sim\_linestore\_tb\_test} and
\code{//docs:dmaw\_vsim\_test}.
\item The assertions are on what the memory holds afterwards, not on
what the beats looked like: every word at its own address, nothing
past the run touched, and a last word holding one byte of the run
over three bytes that were there before. The memory is poisoned with
\code{0xaaaaaa00} first, because without it a correct strobe and an
ignored one both leave zeros there.
\item Sixty-four words in 94 cycles, under the same bound as the read
direction.
\end{itemize}

\subsection*{Limits}

\textbf{It counts the beats of a burst itself, and a write burst is
finished by its response rather than by its last beat.} So a
peripheral that answers before it has taken every beat leaves the
rest of the burst in the channel, and what follows is neither refused
nor written where it was meant to go. Against the Wishbone bridge
before issue 471 nine words of sixty four landed and the rest of
memory was left as it was.

That failure is silent, which is what makes it worth stating here: a
read that stops announces itself by never finishing, and a write that
stops looks exactly like a write that worked.


The byte count's tail is the low two bits of \code{bytes}, so a run
of zero bytes asks for zero beats and does nothing, and nothing
checks that a caller meant that.

The same 4\,KB boundary rule as \code{LineFetch}: a burst must not
cross one, this does not check it, and \code{BEATS} of 256 is
1\,KB, so a 1\,KB aligned base satisfies it without a special case.

It runs on one clock, \code{DefaultClock}. It writes and never reads.
